\documentclass[10pt,a4paper]{article}
\usepackage[spanish,es-noquoting]{babel}
\usepackage[T1]{fontenc}
\usepackage[utf8]{inputenc}
\usepackage{lmodern}
\usepackage{microtype}
\usepackage[a4paper,margin=2.35cm,headheight=14pt]{geometry}
\usepackage{amsmath,amssymb}
\usepackage{siunitx}
\sisetup{output-decimal-marker={,},per-mode=symbol}
\usepackage{circuitikz}
\usepackage{tikz}
\usepackage{booktabs,array,tabularx,longtable}
\usepackage{enumitem}
\usepackage{fancyhdr}
\usepackage{float}
\usepackage{listings}
\usepackage[hidelinks]{hyperref}
\setlength{\parindent}{0pt}
\setlength{\parskip}{0.30em}
\setlist[itemize]{leftmargin=1.35em,itemsep=0.1em,topsep=0.2em}
\setlist[enumerate]{leftmargin=1.5em,itemsep=0.1em,topsep=0.2em}
\renewcommand{\arraystretch}{1.06}
\pagestyle{fancy}
\fancyhf{}
\fancyfoot[C]{\thepage}
\renewcommand{\headrulewidth}{0pt}
\renewcommand{\footrulewidth}{0pt}
\hypersetup{colorlinks=false,pdfborder={0 0 0}}
\lstset{basicstyle=\ttfamily\small,keywordstyle=\color{black},commentstyle=\color{black},stringstyle=\color{black},frame=single,framerule=0.5pt,columns=fullflexible,keepspaces=true,showstringspaces=false,breaklines=true,xleftmargin=0.3em,xrightmargin=0.3em,aboveskip=0.7em,belowskip=0.7em}
\newcommand{\titular}[2]{%
  {\LARGE\bfseries #1\par}
  \vspace{0.45em}
  {\normalsize #2\par}
  \vspace{0.45em}\hrule\vspace{0.9em}
}

\begin{document}
\titular{Detector de paso por infrarrojos}{Registro de pruebas y correcciones \hfill 06/11/2026}

\section{Situación inicial}
Estoy desarrollando un detector de paso para la puerta de una habitación. Mi idea es comprobar primero los componentes en protoboard y, cuando el sistema funcione de forma repetible, preparar una PCB propia en KiCad. Para las pruebas utilizo un ESP32-S3 N16R8, pero no quiero colocar una placa de desarrollo en la versión final.

El montaje actual se divide en dos bloques: un LED IR L-53F3C controlado por un BC337 desde GPIO4 y un fototransistor BPW40 que se lee desde GPIO5 con una resistencia de $\SI{10}{\kilo\ohm}$.

\section{Pruebas realizadas}
\subsection{Salida digital del ESP32-S3}
Primero probé GPIO4 con un LED visible. Utilicé \texttt{pinMode()} para configurar la salida y \texttt{digitalWrite()} para cambiar el estado. Comprobé una tensión próxima a $\SI{3.3}{\volt}$ en HIGH y próxima a $\SI{0}{\volt}$ en LOW.

Al principio el programa cambiaba de estado cada segundo. Para medir con el multímetro era más cómodo aumentar el tiempo o dejar el GPIO fijo en HIGH o LOW. Así pude separar una lectura transitoria del valor estable.

\subsection{Driver del LED infrarrojo}
Monté el BC337 con una resistencia de $\SI{1}{\kilo\ohm}$ en la base y otra de $\SI{180}{\ohm}$ en serie con el LED IR. Las primeras medidas no eran claras y también tuve dudas al diferenciar el emisor y el receptor, que venían juntos en la misma bolsa.

Después repetí las pruebas con el GPIO fijo y anoté:
\begin{center}
\begin{tabularx}{\linewidth}{@{}X r@{}}
\toprule
Medida & Resultado\\
\midrule
GPIO4 con HIGH & $\simeq\SI{3.3}{\volt}$\\
Colector--GND con HIGH fijo & $\SI{0.009}{\volt}$\\
Tensión en R1 ($\SI{180}{\ohm}$) & $\simeq\SI{2.16}{\volt}$\\
Corriente calculada a partir de R1 & $\SI{12.0}{\milli\ampere}$\\
\bottomrule
\end{tabularx}
\end{center}

\begin{equation}
I_{LED}=\frac{2{,}16}{180}=\SI{12.0}{\milli\ampere}.
\end{equation}

La corriente anterior está \emph{calculada a partir de una medida de tensión}; no es una lectura directa de corriente. El valor de colector muy próximo a masa y la caída de tensión en R1 son coherentes con el circuito conduciendo.

\subsection{Receptor BPW40: lectura bloqueada}
Al introducir el fototransistor, el ADC mostraba siempre el valor máximo, 4095, aunque cambiase la iluminación. Revisé el montaje y encontré que había invertido colector y emisor.

Corregí la orientación del componente y comprobé que la salida empezaba a cambiar con la luz. Esta fue la primera avería concreta que pude localizar y corregir.

\subsection{Primeras medidas analógicas}
Con una resistencia pull-up de $\SI{10}{\kilo\ohm}$ obtuve:
\begin{center}
\begin{tabularx}{\linewidth}{@{}X r@{}}
\toprule
Condición & Lectura ADC\\
\midrule
Iluminación normal & 3940--3945\\
Mano interpuesta en la zona del haz & 4057--4067\\
Fototransistor totalmente tapado & 4095\\
\bottomrule
\end{tabularx}
\end{center}

Con nuestro circuito, más luz hace bajar la tensión en GPIO5 y por tanto suele bajar la lectura ADC. Las medidas tienen sentido, pero todavía no puedo atribuir la diferencia solo al LED IR porque el sensor también recibe luz ambiental.

\subsection{Orientación de los componentes}
En una fotografía del montaje vi que el emisor y el receptor estaban verticales en la protoboard, apuntando hacia arriba. Quiero colocar las puntas una frente a otra, con una separación inicial de unos centímetros, y después repetir las medidas. La mejora de orientación todavía no tiene resultados registrados.

\section{Qué he aprendido de estas pruebas}
\begin{itemize}
\item Una medida de tensión se realiza entre dos nodos; para obtener corriente a partir de la resistencia uso $I=V/R$. Si utilizo el multímetro en modo mA, tengo que insertarlo en serie.
\item Si el código cambia más rápido de lo que puedo seguir en la pantalla del multímetro, conviene usar un estado fijo durante el diagnóstico.
\item El fototransistor tiene polaridad y una conexión invertida puede dejar el ADC en el valor máximo.
\item Una lectura distinta al tapar el sensor no demuestra por sí sola que esté detectando exclusivamente la luz del emisor.
\end{itemize}

\section{Estado del prototipo}
\begin{center}
\begin{tabularx}{\linewidth}{@{}X l@{}}
\toprule
Bloque & Estado\\
\midrule
GPIO4 como salida digital & Comprobado\\
BC337 y LED IR & Comprobado eléctricamente\\
BPW40 y lectura ADC & Responde a la luz\\
Alineación óptica & Pendiente de repetir medidas\\
Comparación con IR apagado y encendido & Código preparado; no validado\\
Detección de paso a través de la puerta & Pendiente\\
Notificación Wi-Fi & Pendiente\\
PCB propia & Pendiente\\
\bottomrule
\end{tabularx}
\end{center}

\section{Siguiente sesión}
\begin{enumerate}
\item Colocar el L-53F3C y el BPW40 con las puntas enfrentadas.
\item Medir el ADC con el emisor apagado y encendido, sin mover ningún componente.
\item Interrumpir el haz con cartón opaco y anotar la diferencia de lecturas.
\item Repetir el mismo ensayo con otra iluminación y a mayor distancia.
\end{enumerate}

Con esas medidas podré escoger un umbral basado en datos y plantear las siguientes mejoras del circuito. Por ahora no doy por terminado el detector.

\section{Actualización: validación óptica (19/11/2026)}
Después de orientar el emisor y el receptor, repetí las medidas comparando el ADC con el LED IR apagado y encendido. La diferencia se separó claramente al bloquear el haz. Los intervalos que obtuve fueron:
\begin{center}
\begin{tabularx}{\linewidth}{@{}Xrr@{}}
\toprule
Iluminación & Haz libre & Bloqueado\\
\midrule
Lámpara & 3378--3385 & 323--638\\
Luz ambiental & 3917--3924 & 391--604\\
Oscuridad & 3913--3916 & 319--739\\
\bottomrule
\end{tabularx}
\end{center}
Entre el menor valor libre y el mayor bloqueado de todos los ensayos hay una separación de $3378-739=2639$ cuentas ADC. El montaje ha respondido bien en estas pruebas, pero todavía no he medido falsos disparos ni rapidez de detección. La documentación completa de esta sesión está en \texttt{informe\_validacion\_sensor\_ir.tex}.

También he comprobado que el programa actual tarda aproximadamente \SI{1.6}{\second} por ciclo, porque alterna dos mediciones con pausas y un retardo final. Antes de probar pasos de personas tengo que reducir ese tiempo, establecer la histéresis y verificar el comportamiento con interrupciones breves.

\end{document}